% Generated draft: requires mathematical and source review.
\documentclass[11pt]{article}
\usepackage[margin=25mm]{geometry}
\usepackage{fontspec}
\setmainfont{Latin Modern Roman}
\usepackage{amsmath,amssymb,mathtools,unicode-math}
\setmathfont{Latin Modern Math}
\usepackage{enumitem,needspace,xurl,hyperref}
\setlength{\parindent}{0pt}
\setlength{\parskip}{5pt}
\setlength{\emergencystretch}{3em}
\newcommand{\N}{\mathbb{N}}
\newcommand{\Z}{\mathbb{Z}}
\newcommand{\Q}{\mathbb{Q}}
\newcommand{\R}{\mathbb{R}}
\newcommand{\C}{\mathbb{C}}
\newcommand{\F}{\mathbb{F}}
\newcommand{\A}{\mathbb{A}}
\newcommand{\PP}{\mathbb{P}}
\newcommand{\E}{\mathbb{E}}
\newcommand{\eps}{\varepsilon}
\newcommand{\dd}{\,\mathrm{d}}
\begin{document}
\section*{UnsolvedMath ID: 30003902}
\textbf{Source problem code:} catalogue entry
\subsection{Definitions and mathematical statement}
Let \(S\) be a surface belonging to one of the three classes named in the source: Coble surfaces, Enriques surfaces, or elliptic K3 surfaces. Write \(k\) for its ground field and \(\operatorname{Aut}(S)\) for its automorphism group. The source does not specify \(k\), its characteristic, whether automorphisms are taken over \(k\) or a geometric extension, or the precise conventions for these surface classes. Consequently, the universal quantification below remains subject to unresolved hypotheses.

For a discrete group \(G\) admitting a torsion-free subgroup \(H\) of finite index, its virtual cohomological dimension is
\[
\operatorname{vcd}(G)=\operatorname{cd}_{\Z}(H),
\]
where \(\operatorname{cd}_{\Z}(H)\) is the projective dimension of the trivial \(\Z[H]\)-module \(\Z\). This value is independent of the choice of such \(H\). The source does not state the hypotheses ensuring that this definition applies to \(\operatorname{Aut}(S)\).

Under the usual smooth-generic-fibre convention, a genus-one fibration is a surjective morphism \(f:S\to\PP^1_k\) with connected fibres whose generic fibre \(C_f\), over \(K=k(\PP^1_k)\), is a smooth geometrically integral curve of genus one. It need not have a section. Its Jacobian has generic fibre \(J_f=\operatorname{Pic}^0(C_f)\), an elliptic curve over \(K\). Under the rational-section interpretation of Mordell--Weil rank, one writes
\[
\operatorname{MW\text{-}rk}(f)
=\operatorname{rank}_{\Z}J_f(K)
=\dim_{\Q}\bigl(J_f(K)\otimes_{\Z}\Q\bigr).
\]
The source does not settle whether this rank is taken over the ground field or geometrically, whether a constant-trace quotient is intended, or whether quasi-elliptic fibrations in positive characteristic are included. The displayed definitions describe a conventional interpretation rather than resolve those omissions.

\textbf{Question supplied by the source.} For every Coble surface, Enriques surface, or elliptic K3 surface \(S\) in the intended setting, is
\[
\operatorname{vcd}\bigl(\operatorname{Aut}(S)\bigr)
=\max_{\substack{f:S\to\PP^1\\ f\text{ a genus-one fibration}}}
\operatorname{MW\text{-}rk}(f)?
\]
The maximum is over all genus-one fibrations on \(S\), not merely a chosen elliptic fibration. No convention for an empty indexing set is supplied, and no hypotheses guaranteeing existence and attainment of this maximum are stated. The ground-field, surface-class, fibration, automorphism, and rank conventions must be fixed before this becomes a fully specified universal assertion.
\subsection{Short English statement}
The question asks whether, for every Coble surface, Enriques surface, or elliptic K3 surface in the intended setting, the virtual cohomological dimension of its automorphism group equals the largest Mordell--Weil rank of the Jacobians of its genus-one fibrations. The source leaves several hypotheses and conventions unspecified. Its partially solved label is supplied dated curation, not an independently verified conclusion.
\subsection{Sources}
\begin{itemize}
\item Local UnsolvedMath registry, supplied complete record: permanent ID \(30003902\), source problem code \texttt{catalogue entry}, title \emph{Virtual Cohomological Dimension of Surface Automorphism Groups}, collection \emph{Oberwolfach Reports Open Problems}. No individual proposer is supplied; the proposed year is recorded as 2018.
\item Supplied report citation: \emph{Classical Algebraic Geometry} (2019), \emph{Oberwolfach Reports}, pp.\ 1983--2031, DOI \texttt{10.4171/owr/2018/33}. \url{https://doi.org/10.4171/owr/2018/33}. The supplied background describes these reports as records of mathematical workshops and their open-problem discussions. No problem-specific page or speaker attribution is supplied.
\item Supplied literature source labelled ``OWR report'': \url{https://ems.press/content/serial-article-files/46758?nt=1}.
\item Supplied literature source labelled ``Recent K3 work'': \url{https://repository.kulib.kyoto-u.ac.jp/dspace/bitstream/2433/292398/1/2024-24.pdf}. No author, title, or theorem locator is supplied for this link.
\item Literature assessment checked at \texttt{2026-08-21}: supplied dated curation, not independently verified. The supplied status is ``partially solved''. The assessment states that the Mordell--Weil rank gives one direction, that equality is known for important generic and special surfaces, and that the reverse inequality in the full Coble, Enriques, and elliptic K3 setting remains open. The background qualifies this as dated literature triage derived from the supplied OWR curation and requests independent verification of time-sensitive status and cited claims.
\item The registry labels the statement \texttt{corrected\_verified} and says that mathematical content was preserved while notation, typography, and omitted definitions were repaired. These are supplied editorial assertions, not mathematical verification performed for this draft.
\end{itemize}
\end{document}
